% introduction.tex -- Enterprise Node Orchestrator Paper (RU-V5)

\section{Introduction}

Enterprise blockchain deployments in regulated industries---industrial
traceability, supply chain management, financial record-keeping---face a
fundamental architectural challenge: how to orchestrate the complete
lifecycle of a zero-knowledge validity proof while maintaining continuous
transaction availability and strict data privacy. Unlike public rollup
sequencers where proof generation latency is amortized across thousands of
users, enterprise validium nodes serve a single organization with
deterministic performance requirements and zero tolerance for data
leakage~\cite{ethereum2024zkrollups}.

Production ZK L2 systems such as Polygon zkEVM~\cite{polygon2024architecture,
hermez2023deepdive}, zkSync Era~\cite{zksync2025lifecycle, zksync2025server},
and Scroll~\cite{scroll2023architecture, scroll2025euclid} decompose node
operations into specialized components---sequencers, aggregators, provers,
relayers---connected by internal queues. Recent work on orchestration
frameworks, particularly the push0 system deployed on
Zircuit~\cite{push0_2025orchestration}, demonstrates that event-driven
orchestration with message bus coordination achieves sub-10~ms overhead with
fault-tolerant recovery. However, these architectures target public L2
infrastructure with horizontally scaled prover pools and are
over-provisioned for enterprise validium deployments where a single node
serves one organization.

This paper addresses the gap between public L2 orchestration patterns and
enterprise validium requirements. We present the design of a pipelined,
event-driven state machine orchestrator for the Basis Network enterprise
ZK validium node deployed on Avalanche. The orchestrator integrates five
previously characterized components---Sparse Merkle Tree (RU-V1), ZK state
transition circuit (RU-V2), L1 state commitment contract (RU-V3), crash-safe
batch aggregation (RU-V4), and Shamir-based data availability committee
(RU-V6)---into a single Node.js service that processes the complete
receive-to-settlement cycle.

Our contributions are threefold. First, we define a typed finite state
machine with six states and 17 transitions that supports pipelined execution:
the node continues ingesting transactions while a proof is being generated
or submitted. Second, we provide comprehensive benchmarks demonstrating that
orchestration overhead is 593~ms per batch (0.66\% of the 90-second budget),
with proof generation consuming 64--85\% of end-to-end latency across all
tested configurations. Third, we quantify the pipeline speedup (1.03--1.29$\times$)
as a function of batch size and prover backend, establishing that the
speedup increases with batch size because witness generation time grows
relative to proving time, enabling greater overlap between preparation and
proving phases. These results, together with six formal invariants for
downstream TLA+ specification, provide a complete blueprint for enterprise
ZK validium node implementation.
